% ============== Section 6: Integrated Sensing & Communication (S36-S42) ===
\section{Integrated Sensing \& Communication}

% --- S36 -----------------------------------------------------------------
\begin{frame}{Why ISAC}
  \begin{itemize}\itemsep2pt
    \item One spectrum, one infrastructure, \alert{two functions}:
          communication + radar-style sensing
    \item Drivers: spectrum scarcity; a dense cellular grid is a ready-made
          sensor network; XR/V2X/drones need perception
    \item IMT-2030 lists ISAC as a \emph{defining} scenario --- not an add-on
    \item Economic pitch: sensing ``for free'' on waveforms you already
          transmit
  \end{itemize}
  \note{[S36, 1 min | clock 0:48:30, after Demo 3] Re-entry: from programming
  the channel to reading it. Frame ISAC as the cellular network growing a
  sense of sight: every base station already illuminates the scene with
  wideband, precisely timed signals --- we merely add the receiver processing.
  The regulatory angle (spectrum reuse) is what makes operators care.}
\end{frame}

% --- S37 -----------------------------------------------------------------
\begin{frame}{Radar refresher in one slide}
  \begin{equation*}
    \tau=\frac{2R}{c},\qquad
    f_D=\frac{2v f_c}{c},\qquad
    \Delta R=\frac{c}{2B},\qquad
    \Delta v=\frac{c}{2 f_c\,T_{\mathrm{obs}}}
  \end{equation*}
  \begin{itemize}\itemsep2pt
    \item Bandwidth buys range resolution; \emph{coherent observation time}
          buys velocity resolution
    \item Matched filtering + coherent integration $=$ the SNR workhorses;
          ambiguity function encodes the trade-offs
    \item All classic DSP --- 6G reuses it wholesale, on comm hardware
  \end{itemize}
  \note{[S37, 1.5 min | clock 0:49:30] Calibration slide for the comm-heavy
  audience --- four formulas, no magic. Numbers for our demo: 30 MHz gives
  5 m range resolution; 9 ms of symbols at 28 GHz gives about 5 m/s velocity
  resolution. Say them now; Demo 4 will print exactly these.}
\end{frame}

% --- S38 -----------------------------------------------------------------
\begin{frame}{OFDM as a radar waveform}
  \begin{equation*}
    D[n,m]=\frac{Y[n,m]}{X[n,m]}
    \;\;\xrightarrow[\text{FFT over }m]{\text{IFFT over }n}\;\;
    \text{range--Doppler map}
  \end{equation*}
  \centering
  \resizebox{0.82\textwidth}{!}{% ISAC transceiver: one OFDM frame, comm + sensing branches (slide 38)
\begin{tikzpicture}[node distance=3.5mm and 5mm]
  \node[lab] (d) {data bits};
  \node[blk, right=of d] (mod) {QPSK\\map};
  \node[blk, right=of mod] (ofdm) {OFDM\\mod.\ $X[n,m]$};
  \node[blk, right=of ofdm, minimum width=10mm] (pa) {PA};
  \node[right=5mm of pa] (ant) {\antsym};
  \draw[arr] (d) -- (mod); \draw[arr] (mod) -- (ofdm);
  \draw[arr] (ofdm) -- (pa); \draw[arr] (pa) -- (ant);

  % targets
  \node[lab, align=center] (t1) at ($(ant)+(16mm,9mm)$) {\Large$\ast$\\[-4pt]\scriptsize target 1};
  \node[lab, align=center] (t2) at ($(ant)+(20mm,-7mm)$) {\Large$\ast$\\[-4pt]\scriptsize target 2};
  \draw[arr, corange, thick] ($(ant.east)+(1mm,1mm)$) to[bend left=12] (t1.west);
  \draw[arr, corange, thick] (t1.south west) to[bend left=8] ($(ant.east)+(1mm,-0mm)$);
  \draw[arr, corange, thick] ($(ant.east)+(1mm,-1mm)$) to[bend right=10] (t2.west);
  \draw[arr, corange, thick] (t2.north west) to[bend right=6] ($(ant.east)+(1.5mm,-2mm)$);
  \node[lab, corange] at ($(ant)+(20mm,2mm)$) {echoes};

  % receive path (monostatic, shared ADC)
  \node[blk, below=11mm of pa, minimum width=10mm] (adc) {ADC};
  \draw[arr] (ant.south) |- (adc.east);
  \node[blk, left=of adc] (fft) {OFDM\\demod.\ $Y[n,m]$};
  \draw[arr] (adc) -- (fft);

  % comm branch
  \node[blk, below left=7mm and -2mm of fft, draw=cblue, fill=cblue!8] (eq)
    {equalize\\ + detect};
  \node[lab, left=of eq] (bits) {bits out\\ \scriptsize(comm)};
  \draw[arr] (fft.south) |- (eq.east);
  \draw[arr, cblue] (eq) -- (bits);

  % sensing branch
  \node[blk, left=of fft, draw=corange, fill=corange!8] (div)
    {$D=\dfrac{Y[n,m]}{X[n,m]}$};
  \node[blk, left=of div, draw=corange, fill=corange!8] (fft2)
    {IFFT$_n$ / FFT$_m$};
  \node[lab, left=of fft2, align=center] (rd) {range--Doppler\\map \scriptsize(sensing)};
  \draw[arr] (fft) -- (div);
  \draw[arr, corange] (div) -- (fft2);
  \draw[arr, corange] (fft2) -- (rd);
  \draw[arr, dashed, gray] (ofdm.south) to[bend right=18]
    node[lab, sloped, below]{known TX symbols} (div.north);
\end{tikzpicture}
}
  \note{[S38, 1.5 min | clock 0:51] The Sturm--Wiesbeck trick, walked on the
  diagram: the receiver knows the transmitted grid X (it is the payload!), so
  element-wise division leaves pure channel; an IFFT across subcarriers gives
  delay, an FFT across symbols gives Doppler. Same ADC output feeds the comm
  branch. One frame, two products.}
\end{frame}

% --- S39 -----------------------------------------------------------------
\begin{frame}{Joint waveform design}
  \begin{itemize}\itemsep2pt
    \item Design axes: power/subcarrier allocation, dedicated sensing symbols
          vs full data reuse, beampattern shaping
    \item Inherent friction: comm wants \alert{randomness} (capacity), radar
          wants \alert{determinism} (clean ambiguity function)
    \item Recent theory: the ``random-data radar'' penalty is quantifiable ---
          and often acceptably small
    \item Sensing-assisted communication closes the loop: radar tracks feed
          beam prediction
  \end{itemize}
  \note{[S39, 1 min | clock 0:52:30] The philosophical core of ISAC: our
  cleanest radar waveforms are deterministic chirps; our best comm signals are
  maximum-entropy noise. Joint design is negotiating between them. The
  reverse direction (sensing helps comm --- beam prediction from tracks) may
  land in products first.}
\end{frame}

% --- S40 -----------------------------------------------------------------
\begin{frame}{Fundamental limits: CRB vs rate}
  \begin{equation*}
    \mathrm{var}(\hat\tau)\;\ge\;
    \frac{1}{8\pi^{2}\,\beta_{\mathrm{rms}}^{2}\,\mathrm{SNR}\cdot N_{\mathrm{obs}}}
    \qquad(\beta_{\mathrm{rms}}\!: \text{RMS bandwidth})
  \end{equation*}
  \begin{columns}[T]
    \column{0.5\textwidth}
    \begin{itemize}\itemsep1pt
      \item Delay CRB: RMS bandwidth is king --- edge subcarriers are gold for
            ranging
      \item Rate--CRB region: the capacity--distortion picture of ISAC
      \item Waterfilling for rate $\neq$ waterfilling for sensing
            $\Rightarrow$ principled compromises
    \end{itemize}
    \column{0.48\textwidth}
    \centering
    \includegraphics[width=0.9\textwidth]{fig_rate_crb.png}
  \end{columns}
  \note{[S40, 1 min | clock 0:53:30] One bound to take home: ranging accuracy
  improves with RMS bandwidth, so power at the band edges helps sensing but
  does nothing for capacity under flat channels --- hence the trade-off curve.
  Walk the three dots: 5 percent of resources already buys cm-class CRB at
  these SNRs; the curve is steep at first --- cheap lunches exist.}
\end{frame}

% --- S41 -----------------------------------------------------------------
\begin{frame}{Monostatic, bistatic, and network sensing}
  \begin{columns}[T]
    \column{0.58\textwidth}
    \begin{itemize}\itemsep2pt
      \item \textbf{Monostatic}: full waveform knowledge; but TX--RX
            \alert{self-interference} is brutal (full duplex or quasi-CW
            operation)
      \item \textbf{Bistatic} (BS$\to$UE, BS$\to$BS): no self-interference;
            synchronization and clock offsets bite instead
      \item \textbf{Network/multistatic}: fusion across cells --- geometry
            diversity, coverage; cellular timing infrastructure helps
    \end{itemize}
    \column{0.4\textwidth}
    \centering
    \begin{tikzpicture}[scale=0.8]
      \node[lab] (a) at (0,2.2) {\antsym};
      \node[lab] at (-0.55,2.2) {mono};
      \node[lab] (t) at (1.7,2.25) {\Large$\ast$};
      \draw[arr, corange] (0.25,2.35) to[bend left=14] (1.5,2.4);
      \draw[arr, corange] (1.5,2.15) to[bend left=14] (0.25,2.1);
      \node[lab] (b1) at (0,0.7) {\antsym};
      \node[lab] at (-0.55,0.7) {bi};
      \node[lab] (t2) at (1.1,1.15) {\Large$\ast$};
      \node[lab] (b2) at (2.3,0.6) {\antsym};
      \draw[arr, corange] (0.2,0.9) -- (0.95,1.05);
      \draw[arr, corange] (1.25,1.0) -- (2.1,0.8);
    \end{tikzpicture}
  \end{columns}
  \note{[S41, 1 min | clock 0:54:30] Monostatic is conceptually clean (Demo 4
  assumes it) but hides the hardest RF problem: your own transmit signal is
  120 dB above the echo. Bistatic swaps that for a sync problem --- clocks at
  the two ends drift. Network sensing is where cellular can beat classical
  radar: many viewpoints, one timing backbone.}
\end{frame}

% --- S42 -----------------------------------------------------------------
\begin{frame}{ISAC in 3GPP}
  \begin{itemize}\itemsep2pt
    \item Rel-19: dedicated \alert{channel-model study for ISAC} (targets with
          RCS enter the stochastic model); SA1 use cases feed RAN
    \item Rel-20 (6G study): sensing among the candidate native features;
          ETSI ISG ISAC in parallel
    \item Honest status: pre-normative --- first commercial ``sensing'' will
          likely be positioning-class network features, then micro-Doppler
          applications
  \end{itemize}
  \demoslot{4 (3 min)}{\textbf{Two moving targets on a range--Doppler map} ---
    computed from a frame that is simultaneously carrying QPSK data.
    \hfill\src{matlab/demo4\_isac\_range\_doppler}}
  \note{[S42 + DEMO 4, 1+3 min | clock 0:55:30] Standards status in 30
  seconds, then MATLAB. Demo 4 narration: the frame being processed carries
  random QPSK --- point at the console line confirming data demodulation, then
  the map: two blobs, annotated at 40 m/+15 m/s and 75 m/-25 m/s, within one
  bin of truth. Note the weaker target sits a few dB lower partly due to bin
  straddling (scalloping). If MATLAB fails: appendix A4. Done by ~0:59:30.}
\end{frame}
